\setcounter{section}{2}

\Needspace{5\baselineskip}
\hypertarget{variational-autoencoders-vae-and-vq-vae}{%
\section{Variational Autoencoders (VAE) and VQ-VAE}\label{variational-autoencoders-vae-and-vq-vae}}

\Needspace{5\baselineskip}
\hypertarget{i.-variational-autoencoders-vae}{%
\subsection{I. Variational Autoencoders (VAE)}\label{i.-variational-autoencoders-vae}}

\Needspace{5\baselineskip}
\hypertarget{core-idea-and-architecture}{%
\subsubsection{(1) Core Idea and Architecture}\label{core-idea-and-architecture}}

\Needspace{4\baselineskip}
\begin{enumerate}
\def\labelenumi{\arabic{enumi}.}
\tightlist
\item
  Core idea: from ``deterministic'' to ``probabilistic''
\end{enumerate}

The problem with traditional autoencoders: the encoder directly compresses input data into a fixed, deterministic latent vector. When generating new data, randomly sampling a point in this latent space may decode into meaningless noise, because there is no guarantee that the point lies on the data manifold learned by the model.

The VAE solution: instead of a single vector, the encoder outputs a probability distribution, usually Gaussian, described by two vectors: mean (μ) and variance (σ²). For each input, the model therefore learns a possible region in latent space rather than an exact point.

Depending on the task, we can use only the VAE encoder or combine the encoder and decoder. As noted above, the decoder must sample in probability space, so the output cannot exactly match the input, but its statistical distribution will match the input distribution. The combined encoder and decoder therefore act as an ``imitator,'' producing new data with the same distribution as the input (and thus supporting data augmentation and other uses).

\begin{enumerate}
\def\labelenumi{\arabic{enumi}.}
\setcounter{enumi}{1}
\tightlist
\item
  Key mechanism: how is this distribution learned?
\end{enumerate}

Encoding: given input data x, the encoder outputs the corresponding latent-distribution parameters μ and log(σ²).

Sampling: randomly sample a latent vector z from N(μ, σ²). This crucial step allows gradients to propagate backward through the random operation during training (using the reparameterization trick).

Decoding: feed the sampled latent vector z into the decoder to attempt to reconstruct the original input as x'.

\Needspace{4\baselineskip}
\begin{enumerate}
\def\labelenumi{\arabic{enumi}.}
\setcounter{enumi}{2}
\tightlist
\item
  Loss function: two objectives
\end{enumerate}

\Needspace{5\baselineskip}
The VAE loss has two parts, perfectly reflecting its design:

Reconstruction loss: measures the difference between the decoder's reconstructed data x' and the original input x (such as mean squared error or cross-entropy). This ensures that encoded information can effectively reconstruct the data.

\Needspace{5\baselineskip}
KL divergence loss: measures the difference between the encoder's distribution N(μ, σ²) and the prior (usually the standard normal distribution N(0, 1)). This term regularizes the model:

It forces all encoded distributions to cluster near the origin, making latent space continuous, smooth, and structured.

It prevents the encoder from ``cheating,'' for example by learning a zero-variance distribution for every input (degenerating into deterministic encoding) or encoding different inputs into widely separated regions.

\Needspace{4\baselineskip}
\begin{enumerate}
\def\labelenumi{\arabic{enumi}.}
\setcounter{enumi}{3}
\tightlist
\item
  Advantages and applications
\end{enumerate}

Strong generative capability: because latent space is continuous and smooth, after training, any point sampled from the prior N(0, 1) can be decoded into a plausible, novel sample (such as an image or text).

Meaningful latent space: smooth latent space permits vector arithmetic (such as ``smiling face'' - ``neutral face'' + ``neutral face'' = ``smiling face'').

Applications: image generation, denoising, missing-data completion, and learning the underlying data manifold.

\Needspace{5\baselineskip}
\hypertarget{evidence-lower-bound-elbo}{%
\subsubsection{(2) Evidence Lower Bound (ELBO)}\label{evidence-lower-bound-elbo}}

\Needspace{5\baselineskip}
In probabilistic generative models, observed data \(x\) are usually assumed to be generated by an unknown latent variable \(z\). The central objective is to maximize the marginal likelihood of the observations (that is, the log-evidence):

\[
\log p(x)=\log\int p(x,z)dz=\log\int p(x\mid z)p(z)dz
\]

\Needspace{5\baselineskip}
The difficulties are:

\begin{enumerate}
\def\labelenumi{\arabic{enumi}.}
\tightlist
\item
  Intractable integration: in complex high-dimensional models (such as neural networks), latent-variable space \(z\) is enormous. The integral has no analytical solution and cannot be evaluated by simple numerical computation.
\item
  Intractable posterior: Bayes' theorem gives the true posterior \(p(z\mid x)=\frac{p(x,z)}{p(x)}\). Since the denominator \(p(x)\) (the integral above) cannot be computed, the true posterior \(p(z\mid x)\) is also unknown.
\end{enumerate}

To address this, variational inference introduces a tractable parameterized distribution \(q(z)\) (usually written \(q(z\mid x)\) in VAEs) to approximate the true posterior \(p(z\mid x)\).

ELBO emerges from this approximation: since we cannot directly maximize \(\log p(x)\), we seek a lower bound and maximize it to indirectly maximize \(\log p(x)\).

\Needspace{5\baselineskip}
We start with the Kullback--Leibler (KL) divergence between the approximate distribution \(q(z)\) and true posterior \(p(z\mid x)\). KL divergence measures the difference between two distributions:

\[
D_{KL}(q(z)\parallel p(z\mid x))=\mathbb{E}_{q(z)}\left[\log\frac{q(z)}{p(z\mid x)}\right]
\]

\Needspace{5\baselineskip}
Substitute Bayes' formula \(p(z\mid x)=\frac{p(x,z)}{p(x)}\) and expand:

\[
D_{KL}(q(z)\parallel p(z\mid x))=\mathbb{E}_{q(z)}\left[\log\frac{q(z)p(x)}{p(x,z)}\right]
\]

\Needspace{5\baselineskip}
Apply logarithm rules:

\[
D_{KL}(q(z)\parallel p(z\mid x))=\mathbb{E}_{q(z)}[\log q(z)-\log p(x,z)+\log p(x)]
\]

\Needspace{5\baselineskip}
Since \(\log p(x)\) is a constant for the observed data \(x\) and independent of latent variable \(z\), it can be taken outside the expectation over \(q(z)\):

\[
D_{KL}(q(z)\parallel p(z\mid x))=\mathbb{E}_{q(z)}[\log q(z)-\log p(x,z)]+\log p(x)
\]

\Needspace{5\baselineskip}
Rearranging yields:

\[
\log p(x)=\mathbb{E}_{q(z)}[\log p(x,z)-\log q(z)]+D_{KL}(q(z)\parallel p(z\mid x))
\]

\Needspace{5\baselineskip}
By nonnegativity of KL divergence (\(D_{KL}\ge 0\)):

\[
\log p(x)\ge \mathbb{E}_{q(z)}\left[\log\frac{p(x,z)}{q(z)}\right]
\]

\Needspace{5\baselineskip}
The term on the right is the desired variational lower bound on log-likelihood (ELBO), usually denoted \(\mathcal{L}\):

\[
ELBO=\mathcal{L}=\mathbb{E}_{q(z)}\left[\log\frac{p(x,z)}{q(z)}\right]
\]

Key implication: the equality \(\log p(x)=ELBO+D_{KL}(q(z)\parallel p(z\mid x))\) shows that when optimizing the model (maximizing ELBO), since \(\log p(x)\) is a fixed property of the data, maximizing ELBO is mathematically strictly equivalent to minimizing KL divergence between the approximate posterior \(q(z)\) and the true posterior \(p(z\mid x)\). This explains why optimizing ELBO brings the approximate distribution closer to the true distribution.

\Needspace{5\baselineskip}
This elegant formula splits the complex objective into two highly intuitive terms:

\Needspace{4\baselineskip}
\begin{enumerate}
\def\labelenumi{\arabic{enumi}.}
\tightlist
\item
\Needspace{5\baselineskip}
  Reconstruction term:
\end{enumerate}

\[
\mathbb{E}_{q(z)}[\log p(x \mid z)]
\]

\begin{itemize}
\tightlist
\item
  Mathematical interpretation: the expected log-likelihood of the model generating real data \(x\), given latent variable \(z\) sampled from approximate posterior \(q(z)\).
\item
  Intuition: this term encourages the model to be a good ``decoder.'' Extracted latent features \(z\) must contain enough information to reconstruct the original input \(x\) as perfectly as possible. In actual code, continuous-valued \(x\) usually corresponds to mean squared error (MSE), while binary images correspond to cross-entropy.
\end{itemize}

\Needspace{4\baselineskip}
\begin{enumerate}
\def\labelenumi{\arabic{enumi}.}
\setcounter{enumi}{1}
\tightlist
\item
\Needspace{5\baselineskip}
  Regularization term:
\end{enumerate}

\[
- D_{KL}(q(z)\parallel p(z))
\]

\begin{itemize}
\tightlist
\item
  Mathematical interpretation: the negative KL divergence between approximate posterior \(q(z)\) and prior \(p(z)\) (usually standard normal \(\mathcal{N}(0, I)\)).
\item
  Intuition: this term acts as an ``information bottleneck'' or regularizer. It requires the learned latent distribution \(q(z)\) not to be arbitrary, but to match the predefined prior as closely as possible. Without it, the model might simply memorize every sample (overfitting), leaving latent space disorganized. With it, latent space is regularized, ensuring good interpolation and sampling capabilities for the generative model.
\end{itemize}

\Needspace{5\baselineskip}
\hypertarget{ii.-vector-quantized-variational-autoencoder-vq-vae}{%
\subsection{II. Vector-Quantized Variational Autoencoder (VQ-VAE)}\label{ii.-vector-quantized-variational-autoencoder-vq-vae}}

\Needspace{4\baselineskip}
\begin{enumerate}
\def\labelenumi{\arabic{enumi}.}
\tightlist
\item
  Core idea: from continuous to discrete
\end{enumerate}

A VAE learns a continuous latent space, whereas VQ-VAE learns a discrete latent space, also called a ``codebook.'' What does it contain? It stores K ``prototype feature vectors.'' Each vector row represents a basic visual pattern. For example, row 5 might represent a ``vertical edge,'' row 10 a ``smooth blue texture,'' and row 50 ``a corner of an arc.''

\Needspace{4\baselineskip}
\begin{enumerate}
\def\labelenumi{\arabic{enumi}.}
\setcounter{enumi}{1}
\tightlist
\item
  Key mechanism: vector quantization
\end{enumerate}

The VQ-VAE encoder first maps the input into continuous feature vectors without constructing a VAE-style Gaussian posterior. The model then finds the closest code vector in a predefined codebook and replaces the continuous feature with this discrete code; this is vector quantization. For images, the encoder produces a spatial feature grid. The feature at each location is quantized to the index of its nearest codebook vector, yielding discrete image tokens.

\Needspace{4\baselineskip}
\begin{enumerate}
\def\labelenumi{\arabic{enumi}.}
\setcounter{enumi}{2}
\tightlist
\item
  Advantages and significance of VQ-VAE
\end{enumerate}

Higher-quality discrete representations: discrete tokens often capture inherent, abstract semantic features (such as object parts and textures) more effectively, making them natural for language, reasoning, and related tasks.

Integration with modern generative models: VQ-VAE itself may not achieve the highest reconstruction quality, but its discrete tokens form a bridge between vision and language models. A powerful autoregressive model (such as a Transformer) can learn to model these tokens.

Process: VQ-VAE compresses an image into a discrete token sequence -\textgreater{} a Transformer learns the distribution of these sequences (the ``grammar of images'') -\textgreater{} to generate a new image, the Transformer first generates a plausible token sequence -\textgreater{} the VQ-VAE decoder decodes the tokens into a pixel image.

Representative applications: the first generation of DALL-E and OpenAI Jukebox both used VQ-VAE variants to compress images and audio into discrete representations, then generated them using Transformers.

\FloatBarrier
\Needspace{5\baselineskip}
\hypertarget{references-40}{%
\subsection{References}\label{references-40}}

\vspace{0.25em}
\begin{itemize}
\tightlist
\item
  Kingma, D. P., \& Welling, M. (2014). \href{https://arxiv.org/abs/1312.6114}{Auto-Encoding Variational Bayes}. ICLR 2014.
\item
  van den Oord, A., Vinyals, O., \& Kavukcuoglu, K. (2017). \href{https://arxiv.org/abs/1711.00937}{Neural Discrete Representation Learning}. NeurIPS 2017.
\item
  Ramesh, A., Pavlov, M., Goh, G., et al.~(2021). \href{https://arxiv.org/abs/2102.12092}{Zero-Shot Text-to-Image Generation}. ICML 2021. (DALL-E)
\item
  Dhariwal, P., Jun, H., Payne, C., et al.~(2020). \href{https://arxiv.org/abs/2005.00341}{Jukebox: A Generative Model for Music}. arXiv:2005.00341.
\end{itemize}

\clearpage
